% Documentclass Settings
	\documentclass[a4paper, twoside]{article}
	% Margins
	\usepackage{geometry}
		\geometry{
			% showframe,
			left   = 20mm,
			right  = 20mm,
			top    = 20mm,
			bottom = 20mm
		}

	% Headers and footers
	\usepackage{fancyhdr}
		\pagestyle{fancy}
		\fancyhf{}
		\fancyhead[LE, RO]{\thepage}
		\fancyhead[CE]{\textsc{Cong Wen}}
		\fancyhead[CO]{\textsc\rightmark}

		\setlength{\headheight}{15pt}
		\renewcommand\headrulewidth{0pt}








\input{../../universal-pre}

\title{\tbf{Counting Lines of a Quintic 3-fold}}
\author{{Cong Wen, 06/2019}\\\tbf{Advisor}: James Fullwood}
\date{}
% Last Updated: July 11, 2022


\begin{document}
\maketitle
\thispagestyle{empty}
\begin{abstract}
    This short note\footnote{Recompiled version on March 06, 2023} aims at illustrating the basic ideas in enumerative geometry and foundations of intersection theory with a basic example: counting lines in a general quintic 3-fold $X\subset\bP^4$, the computation played an important role in the discovery of mirror symmetry. The basic idea of enumerative geometry is to count points in parameter space which satisfies given conditions rather than literally counting lines. Since the 3-fold $X$ is contained in $\bP^4$, we first consider the Grassmannian $G(2,5)$ consisting of all the lines in $\bP^4$. Then we notice the correspondence between lines on $X$ and the top Chern class of the section of a vector bundle over $G$. Making use of the splitting principle, we can reduce the problem to calculating Chern class on a line bundle of $G$, which can be further dealt with by Schubert cycles.

\end{abstract}
\tableofcontents


% ----------------------------------------------------------------------------------------------------

\sct{Introduction}
Intersection theory is at the heart of algebraic geometry, at the beginning, the fact that the number of solutions to a system of polynomials equations is constant in many circumstances as we vary the coefficients has fascinated geometers. B\'ezout's theorem is a first example, which stated that the total number of intersection points of two complex projective curves counted with their multiplicities, is equal to the product of their degrees.

At the beginning of $19$th century, algebraic geometers has made two important choices: work over complex numbers rather than real numbers, and work in projective spaces rather than affine spaces. And the introduction of parameter space is introduced when counting solutions. Despite a lack of precise foundations, Schubert has pushed enumerative geometry to impressive heights, who calculated the number of twisted cubics tangent to $12$ quadrics ($5819539783680$), and to put his method---schubert cycles on a rigorous foundation is the well know Hilbert's fifteenth problem. At the end of $20$th century, Fulton's landmark book \tit{Intersection Theory}\cite{Fulton1998} for the first time put the subject on a precise and rigorous foundation. The quest for a rigorous foundation not only drove much of algebraic geometry, but also took an important role in the discovery of mirror symmetry.

In this short article we will not present proofs of every theorem, instead, we aims at illustrating how powerful the intersection theory is when solving classical enumerative problems by using Chern class to solve a concrete counting problem---how many lines are contained in a quintic $3$-fold?

\sct{Definitions}
\ssc{The Motivating Problem}
\begin{dfn}
	A complex projective space $\bP^n$ is defined as the quotient of $\bC^{n+1}/C^*$, in which $C^*$ acts on $\bC^{n+1}$ by linear multiplication, $\lambda\cdot(a_0,\cdots,a_n)\equiv (\lambda a_0,\cdots, \lambda a_n)$. And By a line in projective space we mean the projectivization of a $2$-dim linear subspace of $\bC^n$.
\end{dfn}

\begin{prb}\label{count}
Let $X\subset\bP^4$ be a hypersurface defined by a degree $5$ homogeneous polynomial (this is a quintic $3$-fold), how many lines are contained in $X$? 
\end{prb}

\ssc{Chow groups}

$X$ is a variety, and given two closed subvarieties $V,W$ in $X$, there are two problems when we talk about the intersection of $V,W$ in general case.

On the one hand, $V\cap W$ might be non-reduced, i.e., they might be tangent to each other at a point of intersection, so we have to count the "multiplicity".

On the other hand, the might not intersect transversely. we say two subvarieties $A,B$ of a variety $X$ intersect transversely at a point $p$ if $A,B,X$ are all smooth at $p$ and the tangent spaces to $A$ and $B$ span the tangent space to $X$. $A,B$ are generically transverse if they meet transversely at a general point of each component $C$ of $A\cap B$.

The introduction of Cycles and Chow group is a solution to them.

Let $X$ be any algebraic variety (or a scheme), then the free abelian group $Z(X)$ generated by the set of subvarieties of $X$ is graded by the dimension of subvarieties. Write $Z=\bigoplus_kZ_k(X)$, and every cycle $Z=\sum n_iY_i$. A divisor is an $n-1$-cycle on a pure $n$-dimensional scheme. If $Y$ is a closed subvariety of $X$ and a subscheme then denote by $\langle Y\rangle$ the formal combination $\sum l_iY_i$, in which $Y_i$ are irreducible components of $Y_{red}$. The coefficient involved is the multiplicity of scheme $Y$ along $Y_i$. And two cycles $A\sum n_iA_i$ and $B=\sum m_jB_j$ are generically transverse if each $A_i$ and $B_j$ are generically transverse.

 Chow group is a algebro-geometric analog of the homology of a topological space. In the case of a smooth variety, the intersection product makes it into a graded ring. The Chow group is the defined as the cycles of $X$ modulo rational equivalence, $A_0,A_1$ are rationally equivalent if they can "rationally varies" to each other. To give a formal definition,
\begin{dfn}
	$\on{Rat}(X)\subset Z(X)$ is the free Abelian group generated by cycles of the form
	$\langle\Phi\cap(\{t_0\}\times X)\rangle-\langle\Phi\cap(\{t_1\}\times X)\rangle$, in which $t_i\in\bP^1$ and $\Phi$ is a subvariety of $\bP^1\times X$ not contained in any $\{t\}\times X$. And two cycles are equivalent if their difference is in $\on{Rat}(X)$.
\end{dfn}

\begin{dfn}
	The Chow group of $X$ is the quotient
	$$
	A(X)=Z(X)/\on{Rat}(X)
	$$
\end{dfn}
In fact, $A(X)$ is graded by dimension $A(X)=\bigoplus A_k(X)$, and $A_k(X)$ is the group of rational equivalence classes of $k$-cycles. The definition means we consider rational equivalence class of a cycle instead of itself, this makes them more flexible and can move rationally. This flexibility made it possible to guarantee the transverse intersection of $V$ and $W$ in their equivalence classes, we see in the lemma,
\begin{lem}[Moving Lemma]
	$X$ is a smooth quasi-projective variety.
	\begin{itm}
		\item Every $\alpha,\beta\in A(X)$ there are generically transverse cycles $A,B\in Z(X)$ with $[A]=\alpha$ and $[B]=\beta$.
		\item The class $[A\cap B]$ is independent of the choice of such $A$ and $B$.
	\end{itm}
\end{lem}

But now the relation of $[A\cap B]$ and $[A],[B]$ is not clear, and it turns out that the ring structure on Chow groups made it computable.
\begin{thm}
	$X$ is a smooth quasi-projective variety, then there is a unique product structure on $A(X)$ satisfying: If two subvarieties $A,B\subset X$ are generically transverse, then
	$$
	[A][B]=[A\cap B].
	$$
	The structure $A(X)$ into a commutative graded ring,
	$$
	A(X)=\bigoplus_{k=0}^{\dim X}A^k(X).
	$$
\end{thm}

\begin{exg}
	Let $H$ be any hyperplane in $\bP^n$, then
	$$
	A(\bP^n) = \mbb{Z}[H]/H^{n+1},
	$$
	and $A_k(\bP^n)$ is the degree $n-k$ monomials.
\end{exg}
This means the computation of intersection number in $\bP^n$ reduces to polynomial calculation. And counting number of points (codimension $n$) of an intersection in $\bP^n$ reduces to find the degree of $H^n$ of its rational equivalence classes in the Chow ring.

We see an easy proof of B\'ezout's theorem from this: two complex projective plane curves $A,B$ of degree $m,n$, there rational equivalence class are $nH,mH\in A_1(\bP^2)$ respectively, so $[A\cap B]=[A][B]=mn H^2$, hence the intersection number is $mn$ when $A,B$ are in general position.

\ssc{Chern class}
Let $\cE$ be a vector bundle on a variety $X$ of dimension $n$, then this section will introduce Chern class $c_k(\cE)\in A_{n-k}(X)$ and its properties. Before that we define the degeneracy locus. First recall the first Chern class $c_1(\cL)$ of a line bundle $\cL$ on a variety $X$:
$$
c_1(\cL)=[\Div(\tau)]\in A_{n-1}(X),
$$
in which $\tau$ is a rational section of $\cL$, in this case we define $c_i(L)=0, i\gqs 2$.

Chern class is a characteristic class describing the nontriviality of a bundle. It is well known in topology that a rank $r$ bundle is trivial if and only if it has $r$ everywhere-independent global sections. More generally, consider $r-i+1$ global sections, and their degeneracy locus is defined as the vanishing of
$$
D=V(\tau_0\wedge\cdots\wedge\tau_{r-i})\subset X.
$$
$D$ is indeed a scheme, and has good properties:
\begin{lem}
	Let $\cE$ be a vector bundle of rank $r$, $\tau_0\cdots,\tau_{r-i}$ be global sections of $\cE$, and $D$ is defined as before, then
	\begin{itemize}
		\item No component of $D$ has codimension $>i$
		\item If $\tau_i$ are general elements of a vector space $W\subset H^0(\cE)$ of global sections generating $\cE$, then $D$ is generically reduced and has codimension exactly $i$ in $X$.
	\end{itemize}
\end{lem}

\begin{cor}\label{locus}
If there is only one global section $\tau_0$ satisfying the preceding conditions, we know the components of $D(\tau_0)$ (which is the zero locus of $\tau_0$) are precisely of codimension $r$, i.e. they are isolated points.
\end{cor}

Now we can characterize Chern class for $\cE$
\begin{thm}
	There is a unique way of assigning to each vector bundle $\cE$ on a smooth quasi-projective variety $X$ a class $c(\cE)=1+c_1(\cE)+\cdots\in A(X)$ in such a way that:
	\begin{itemize}
		\item If $\cL$ is a line bundle on $X$ then the Chern class of $\cL$ is $1+c_1(\cL)$, in which $c_1(\cL)$ is the class of divisor $D$ of any rational section $\tau$ of $\cL$
		\item If $\tau_0\cdots,\tau_{r-i}$ are global sections of $\cE$, and the degeneracy locus $D$ has codimension $i$, then $c_i(\cE)\in A^i(X)$
		\item (Whitney's Formula) If $0\rightarrow\cE\rightarrow\cF\rightarrow\cG\rightarrow 0$ is a short exact sequence of vector bundles on $X$, then $c(\cF)=c(\cE)c(\cG)$.
		\item (Functoriality) If $\varphi:Y\rightarrow X$ is a morphism of smooth varieties, then $\varphi^*(c(\cE))=c(\varphi^*(\cE))$.
	\end{itemize}
\end{thm}

Chern class has another remarkable result which can be powerful when using together with Whitney's formula:
\begin{thm}[Splitting Principle]
	Any identity among Chern classes of bundles that is true for bundles that are direct sum of line bundles is true in general.
\end{thm}

If $\cE$ is a vector bundle of rank $r$
\begin{exg}
	$c_i(\cE)=0$ for $i>r$.
\end{exg}
The reason is easy using the splitting principle, since if $\cE$ splits into $\oplus_i^r\cL_i$, then $c(\cE)=\prod_i^r(1+c_1(\cL_i))$ has no term of degree $>r$.
\begin{exg}
	$c_i(\cE^*)=(-1)^ic_i(\cE)$
\end{exg}
The proof is as easy as the preceding one, use $c_1(\cL^*)=-c_1(\cL)$ when $\cL$ is a line bundle.

\sct{Grassmannians}\label{grass}
This part we will describe the structure of Chow ring of Grassmannian, in particular that for $G(2,5)$, this is a fundamental part in the following calculation.

In fact, Grassmannian is a generalization of projective plane, and we can show that $G(2,n+1)$ is in one-one correspondence to all the lines in $\bP^n$, we only illustrate the method by analyze the case $G(2,5)$, which plays an important role in computing lines in a quintic 3-fold, as for the general case $G(k,n)$, we recommend interested readers to book by Harris.

\ssc{Projectivization}
For any linear $V$ over $\bC$, we can take the projectivization of $V$ by identifying $v\sim\lambda v$, in which $v$ is a non zero vector in $V$ and $\lambda$ is a non zero complex number, denote by $P(V)$ the resulting space. If $\dim V=n<\infty$, then $V\cong\bC^n$, so we can consider $V$ as $\bC^n$.

Recall that the original definition of a Grassmannian $G(k,n)$ is all the $k$-subspaces of a vector space of dimension $n$,
$$
G(k,n)=\{W:W\subset\bC^n, \dim W=k\},
$$

So it is easy to see $G(1,n+1)$ is exactly the projective $n$-space given a topological structure, also we have $G(1,n+1)\cong P(\bC^{n+1})=\bP^n$.

Although $\bP^n$ is not a vector space, but we can also define its lines and $n$-subspace via the projectivization of corresponding linear subspace of one more dimension in $\bC^n$, then all the lines in $\bP^n$ can be written as
$$
\{P(W)\subset\bP^n:W\subset\bC^{n+1},\dim W=2\}.
$$

By the abuse of notation we just denote by $G(k+1,n+1)$ all the projective $k$-planes in $\bP^n$ without unambiguity, for example, $G(2,n+1)$ all the lines in $\bP^n$. We restrict out attention to $G(2,5)$ in the following context and write it as $G$ for simplicity.

\ssc{Cell Decomposition}
To describe the structure and compute homology groups of $G(2,5)$, we may consider the cellular decomposition of $G(2,5)$. Before introducing Schubert cycle notations, we first consider its structure as a smooth manifold.

Any $k$-dimensional space can be written as a $2\times 5$ matrix $\Lambda$  of rank $2$.
$$
\mat{v_{11}&\cdots& v_{15}\\v_{21}&\cdots&v_{25}}
$$

Then two matrix $\Lambda_1,\Lambda_2$ are the same $2$-dimensional space iff for some $G\in\on{GL}_2(V)$, $G\Lambda_1=\Lambda_2$. Take a multiindex $I=\{i_1,i_2\}\subset\{1,\cdots,5\}$, let $V_{I^\circ}=\langle e_i:i\notin I\rangle\subset\bC^5$ and $U_I=\{\Lambda\in G(2,5):\Lambda\cap V_{I^{\circ}}=\emptyset\}$. So $U_I$ contains precisely those $\Lambda\in G(2,5)$ whose $I$th minor of the matrix representation is a nonsingular $2\times 2$ matrix.

So we can choose a unique matrix representation of $\Lambda$ to be the one with its $I$th minor being an identity matrix. For example, the matrix representation of $U_{\{1,2\}}$ is
$$
\left\{\mat{1&0&v_{13}&v_{14}&v_{15}\\0&1&v_{23}&v_{24}&v_{25}}:v_{ij}\in\bC\right\}\cong\bC^6.
$$

Also, write the isomorphism as $\varphi_I$, then all the possible $I$ with corresponding $U_I$ forms manifold structure, in addition, $\varphi_I\circ\varphi_{I'}^{-1}$ is holomorphic, to see this, consider $\Lambda\in U_I\cap U_{I'}$, and denote $\Lambda^I$ and $\Lambda^{I'}$ the matrix representation respectively. Then if $\Lambda_{I'}^I$ is the $I'$th minor of $\lambda^{I}$, we have $\Lambda^{I'}=(\Lambda_{I'}^I)^{-1}\Lambda^I$, so $\varphi_I\circ\varphi_{I'}^{-1}$ is a linear transformation for every line in $G(2,5)$, and since the entries vary holomorphically, $\varphi_I\circ\varphi_{I'}^{-1}$ is holomorphic. So $G$ is a smooth manifold.

Since every element in $G(2,n+1)$ can be characterized by a $2$-dimensional subspace in $\bC^{n+1}$, so we first define subspaces $G_{i,j}\subset\bC^{n+1}$ and assume $0\lqs j\lqs i\lqs n-1$,
$$
G_{i,j}:=\mat{a_1&\cdots&a_{n-1-i}&1\\b_1&\cdots&b_{n-1-i}&0&b_{n-i}&\cdots&b_{n-1-j}&1}\cong\bC^{2n-2-i-j}.
$$

then $\{G_{i,j}\}_{0\lqs j\lqs i\lqs n-1}$ is a family of cells contained in $G(2,n+1)$ and the largest one is $G_{0,0}$ of dimension $2n-2$, in fact, by change of coordinates, all corresponding $G_{0,0}$ form a cover of $G$. To see this, consider that all the lines which do not intersect with the algebraic set $Z(x_{n-1}, x_n)$ are in $G_{0,0}$ and vice versa, so for any line $l=P(\langle e_1,e_2\rangle), e_i\in\bC^{n+1}$ (view $e_i$ as linear transformation of $x_i$), we have $l\cap Z(e_1, e_2)=\emptyset$ hence $l\in G_{0,0}$ (under the coordinate $\{e_i\}$).

Further, on the intersection of covers, transition function is a linear transformation, hence $G$ is a smooth manifold.

For a formal definition of \emph{Schubert cycle}\footnote{the rigorous footing of Schubert calculus is Hilbert's fifteenth problem, and was well established by Fulton in his book, intersection theory}, first note that obviously the choice of basis should not matter, and denote by $\{e_i\}_{i=1}^n$ the basis of $\bC^{n+1}$, and let $V_i=\langle e_1,\cdots, e_i\rangle$, so we have the \emph{flag} $0=V_0\subset\cdots\subset V_{n+1}=\bC^{n+1}$, we define $\sigma_{a,b}$ as follows.\footnote{in the case of $G(k,n)$, the Schubert cycle is defined similarly using the multi index $a=(a_1,\cdots,a_k)$, and $\sigma_{a}=\{\Lambda\in G(k,n):\dim(\Lambda\cap V_{n-k+i-a_i})\gqs i\}$}
$$
\sigma_{i,j}=\{l\in G: l\cap P(V_{n-i})\neq\emptyset, l\subset P(V_{n+1-j}) \}.
$$

In fact, we can check that $\oL{G_{i,j}}=\sigma_{i,j}$, this is exactly the geometric interpretation of $G_{i,j}$. To see this, first note that $G_{i,j}\subset\sigma_{i,j}$, further imagine an element defined by the matrix representation and $a_1,\cdots a_{n-1-i}$ tends to infinity with a given ratio, this will gives an element exactly in $G_{i-1,j}$ or $G_{i,j-1}$ if they exist, but this does not have an effect on satisfying condition of $\sigma_{i,j}$.

Also, we have $\sigma_{a,b}\in H^{2a+2b}(G)$, the integral cohomology of $G$, and the Pieri's formula to do Schubert calculus in practice,
\begin{thm}[Pieri's Formula]
	For any Schubert class $\sigma_b,c\in A(G)$ and any integer $a$,
	$$
	\sigma_{a,0}\cdot\sigma_{b,c}=\sum_{d+e=a+b+c, d\gqs b\gqs e\gqs c}\sigma_{d,e}.
	$$
\end{thm}

Using this formula, we can compute the whole multiplication table, and further more, the cellular decomposition itself gives the structure of Chow ring of $A(G)=\oplus_i A_i(G)$, for instance $G=G(2,5)$, the multiplication table is also shown below.
\begin{align*}
	A_6(G)&=\mbb{Z}\cdot 1,\\
	A_5(G)&=\mbb{Z}\cdot\sigma_{1,0},\\
	A_4(G)&=\mbb{Z}\cdot\sigma_{2,0}\oplus\mbb{Z}\cdot\sigma_{1,1},\\
	A_3(G)&=\mbb{Z}\cdot\sigma_{2,1}\oplus\mbb{Z}\cdot\sigma_{3,0},\\
	A_2(G)&=\mbb{Z}\cdot\sigma_{2,2}\oplus\mbb{Z}\cdot\sigma_{3,1},\\
	A_1(G)&=\mbb{Z}\cdot\sigma_{3,2},\\
	A_0(G)&=\mbb{Z}\cdot\sigma_{3,3},\\
\end{align*}
\begin{table}[H]
	\centering
	\caption{Multiplication Table of A(G)}
	\label{tab:multi}
    \begin{tabular}{c|cccccccccc}
    	\toprule
      	 & 1 & $\sigma_{1,0}$ & $\sigma_{2,0}$ & $\sigma_{3,0}$ & $\sigma_{1,1}$ & $\sigma_{2,1}$ & $\sigma_{2,2}$ & $\sigma_{3,1}$ & $\sigma_{3,2}$ & $\sigma_{3,3}$ \\
     	\midrule
     	1 & 1&$\sigma_{1,0}$ & $\sigma_{2,0}$ & $\sigma_{3,0}$ & $\sigma_{1,1}$ & $\sigma_{2,1}$ & $\sigma_{2,2}$ & $\sigma_{3,1}$ & $\sigma_{3,2}$ & $\sigma_{3,3}$ \\
     	$\sigma_{1,0}$ &$\sigma_{1,0}$&$\sigma_{2,0}$+$\sigma_{1,1}$&$\sigma_{3,0}$+$\sigma_{2,1}$ &$\sigma_{3,1}$&$\sigma_{2,1}$&$\sigma_{3,1}$+$\sigma_{2,2}$&$\sigma_{3,2}$&$\sigma_{3,2}$&$\sigma_{3,3}$\\
     	$\sigma_{2,0}$ &$\sigma_{2,0}$&$\sigma_{3,0}$+$\sigma_{2,1}$&$\sigma_{3,1}$+$\sigma_{2,2}$&$\sigma_{3,2}$&$\sigma_{3,1}$&$\sigma_{3,2}$&&$\sigma_{3,3}$\\
     	$\sigma_{3,0}$ &$\sigma_{3,0}$&$\sigma_{3,1}$&$\sigma_{3,2}$&$\sigma_{3,3}$\\
     	$\sigma_{1,1}$ &$\sigma_{1,1}$&$\sigma_{2,1}$&$\sigma_{3,1}$&&$\sigma_{2,2}$&$\sigma_{3,2}$&$\sigma_{3,3}$\\
     	$\sigma_{2,1}$ &$\sigma_{2,1}$&$\sigma_{3,1}$+$\sigma_{2,2}$&$\sigma_{3,2}$&&$\sigma_{3,2}$&$\sigma_{3,3}$\\
     	$\sigma_{2,2}$ &$\sigma_{2,2}$&$\sigma_{3,2}$&&&$\sigma_{3,3}$\\
     	$\sigma_{3,1}$ &$\sigma_{3,1}$&$\sigma_{3,2}$&$\sigma_{3,3}$\\
     	$\sigma_{3,2}$ &$\sigma_{3,2}$&$\sigma_{3,3}$\\
     	$\sigma_{3,3}$ &$\sigma_{3,3}$\\
     	\bottomrule
     \end{tabular}
\end{table}


%\subsection{Line Bundles on $G(2,n+1)$}
%We now introduce the tautological line bundle, and the Chern class of $G=G(k,n)$, we want to give a description of the tautological bundle $E$ and its dual bundle $Q$ on $G$ by give a concrete expression of their transition functions.
%
%The tautological bundle $\pi:E\rightarrow G$ of Grassmannian $G$ is a bundle whose fiber attached to every $k$-dimensional space $X\subset\bC^n$ is itself, namely, $\pi^{-1}(X)=X\times X$. To illustrate the underlying charts and how they are glued together, we denote by $p_{XY}$ the orthogonal projection map between two $k$-dimensional space, $p_{XY}:X\rightarrow Y\subset\bC^n$, obviously $p_{XY}\in\\on{End}(\bC^n)$ is of rank $k$.
%
%Given $X\in G$, let $V=\{Y\in G(k,n):\dim p_{YX}(X)=k\}$ be a neighborhood of $X$, note that every $Y\in V$, $p_{YX}(Y)\cong X$.
%
%The local trivialization in a neighborhood of $X$ is then given by
%\begin{align*}
%	\phi:\pi^{-1}(V)&\rightarrow G(k,n)\times X\\
%	(Y,v)&\mapsto(Y,p_{YX}(v)).
%\end{align*}
%
%In fact, if $X=X_I=\langle e_i:i\notin I\rangle$ for some multiindex $I\subset\{1,\cdots,n\}$, then $V$ defined above is precisely $U_I$ in the previous part, a chart of the manifold structure of $G$, and we relabel $\phi$ as $\phi_I$. Hence $\{\pi^{-1}(U_I)\}_{I\subset\{1,\cdots,n\}}$ forms an open cover of the tautological bundle of $G$, the transition function $g_{IJ}$ can be given by the trivialization function: If $X\in U_I\cap U_J$, then for all $(X,v)\in\pi^{-1}(U_I)$, $(X,g_{IJ}(X)v)=\phi_J\circ\phi_I^{-1}(X,v)$, since $v$ is arbitratily chosen, $g_{IJ}(X)$ is uniquely determined.
%
%Then we consider the dual bundle $\pi^*:Q\rightarrow G$ of $E$, whose fiber $Q_X$ is the dual space of $E_X$, all the linear functionals on a $k$-dimensional space $X$, $Q$ is given the same charts as $E$,  the transition function is given by $g_{ij}^*(X)=(g_{ij}^{H}(X))^{-1}$. To see this, we naturally require that the inner product is well defined under the action of transition function, namely, if $X\in U_I\cap U_J$ and $(X,u)\in\pi^{-1}(U_I), (X,v)\in(\pi^*)^{-1}(U_I)$, in each bundle, they are equivalent to $(X,g_{IJ}(X)u)$ and $(X,g_{IJ}^*(X)v)$ respectively (write $g_{IJ}(X)$ as $g$ in brief), so we naturally require $v(u)=(g^*v)(gu)$, then $\langle v,u\rangle=\langle g^*v,gu\rangle=\langle g^Hg^*v,u\rangle$, since this holds for all $u$ and $v$, we have $g^*=(g^H)^{-1}$.


%To give a concrete example, take $G=G(2,3)$, and the field be $\mbb{R}$. Then the charts are $\{Zi\}$, in which $Z_i=Z(x_i), i=1,2,3$. Let $X$ be the plane perpendicular to $(-2,1,0)$, then $X\in U_1\cap U_2, X\notin U_3$. We now calculate $g_{12}(X)$ and $g_{12}^*(X)$. Write $X=\langle\frac{1}{\sqrt{5}}\vec{i}+\frac{2}{\sqrt{5}}\vec{j},\vec{k}\rangle=E_X$, then $(X, (a,b)^T)\in E$ is trivialized as $(X_1, (\frac{2}{\sqrt{5}}a, b)^T)$ in $U_1$ and $(X_2, (\frac{1}{\sqrt{5}}a, b)^T)$ in $U_2$, so $g_{12}(X)=\mat{1/2&0\\0&1}$. In this case, we naturally need the inner product is invariant under the transition map. If $(X,h)\in(X,Q_X), Q_X=X ^*$ and globally $h=x\vec{e}_i+y\vec{e}_j+z\vec{e}_k$, then
%$$
%\frac{2}{\sqrt{5}}ay+bz=\frac{1}{\sqrt{5}}ax+bz
%$$
%Hence we have $x=2y$, in the open set $U_1$ and $U_2$, $h$ is clearly represented as $(y,z)^T$ and $(x,z)^T$, so $g_{12}^*(X)=\mat{2&0\\0&1}=(g_{12}(X)^T)^{-1}$.

\sct{Solution to the Problem}

\ssc{Strategy}
Up to now we can state the strategy of the solution to Prob. \ref{count} raised at beginning. We first pack up all the lines in $\bP^4$ into a parameter space---this is the Grassmannian $G(2,5)$ shown above. A crucial observation shows that the condition $l$ is in the hypersurface $X$ can be restated as the vanishing of a global section of a vector bundle on $G(2,5)$. Then counting lines in a quint $3$-fold can be shown to equivalent to find out the number of zero locus of it, which is the degree of top Chern class by Cor. \ref{locus}. To complete the final the concrete computation, we are going to use Splitting principle together with Whitney's formula, since now Chern classes live in the Chow ring of $G(2,5)$, we need the multiplication structure of it, the multiplication table is given in Sec. \ref{grass}.

\ssc{Calculation}
We first show how the counting number of lines reduce to couting zero locus of a global section of a vector bundle, and perform the concrete calculation below.

Consider the 3-fold $X\subset\bP^4$ which is defined by a homogeneous polynomial $F$ of degree $5$, which can be regarded as a vector, namely $\mbf{v}_F$, in the space $V=\langle x_0^5, x_0^4x_1,\cdots, x_4^5\rangle$ whose dimension is $126$.

Then consider the vector space $Q_l$, spanned by linear polynomials on a line $l$, then the quintic homogeneous polynomial on this line should be the fifth symmetric power of $Q_l$ and $\dim\Sym^5Q_l=6$.

A crucial observation is that if $Q_l=\langle e_1,e_2\rangle$, in which $e_1, e_2$ are linear combination of $x_0,\cdots,x_4$, and the homogeneous coordinates of $l$ in terms of $Q_l$ is $[k_1,k_2]\in\bP^1$. By change of homogeneous coordinates, $F$ can be rewritten as a polynomial $\tilde{F}$ in terms of $\Sym Q_l$, hence $F|_l=0$ implies $\tilde{F}$ vanishes for all $[k_1,k_2]\in\bP^1$, so the coefficients must be zero. Thus we conclude that the projection of the vector $\mbf{v}_F$ on the subspace $\Sym^5Q_l$ must be zero if $F|_l=0$. So every $l\subset X$ corresponds to a subspace $Q_l$ on which $\mbf{v}_F$ has zero projection.

Then as $l$ varies in the Grassmannian $G$, $Q_l$ varies correspondingly to form a vector bundle $Q$ of rank $6$ over $G$ (this is the dual bundle of tautological bundle of $G(2,5)$), and the given polynomial $F$ determines a global section $\pi$ of $\Sym^5Q$ by projecting $\mbf{v}_F$ to $\Sym^5Q_l$,
\begin{align*}
	\pi:G&\rightarrow\Sym^5Q\\
	l&\mapsto\mbf{v}_F|_{\Sym^5Q_l}
\end{align*}

And the zero locus of $\pi$, which corresponds to the top Chern class of $\Sym^5Q_l$ exactly corresponds to lines which belong to $X$. To count the number of zero locus of $\pi$, by Cor. \ref{locus}, we only need to find the degree of top Chern class of $\Sym^5Q$.
$$
n=\on{deg}c_6(\Sym^5Q).
$$
 
By splitting principle, assume $Q$ is decomposed into direct sum of line bundles $Q\cong\cE\oplus\cF$ we can compute the Chern class of $Q$, denote by $\alpha,\beta$ the first Chern class of $\cE,\cF$ respectively,
$$
c(Q)=1+(\alpha+\beta)+\alpha\beta
$$
But we already know $c(Q)=1+\sigma_{1,0}+\sigma_{1,1}$ (the proof is given in\cite{GH1994}), this shows that
$$
\alpha+\beta=\sigma_{1,0}, \alpha\beta=\sigma_{1,1}.
$$

Now the Chern class of $\Sym^5Q$ can be computed, by Whitney's formula,
\begin{equation*}
	\begin{split}
		\Sym^5Q&\cong\cE^{\otimes 5}\oplus (\cE^{\otimes 4}\otimes \cF)\oplus (\cE^{\otimes 3}\otimes \cF^{\otimes 2})\oplus (\cE^{\otimes 2}\otimes \cF^{\otimes 3})\oplus (\cE^{\otimes 1}\otimes \cF^{\otimes 4})\oplus \cF^{\otimes 5}\\
		c(\Sym^5Q)&=(1+5\alpha)(1+4\alpha+\beta)(1+3\alpha+2\beta)(1+2\alpha+3\beta)(1+\alpha+4\beta)(1+5\beta)\\
%		&=600 \alpha^5\beta + 3850 \alpha^4 \beta^2 + 6725 \alpha^3 \beta^3 + 3850 \alpha^2 \beta^4 + 600 \alpha \beta^5\\
%		&+ 120 \alpha^5  + 2140 \alpha^4 \beta + 7115 \alpha^3 \beta^2 + 7115 \alpha^2 \beta^3 + 2140 \alpha \beta^4 + 120 \beta^5\\
%		&+ 274 \alpha^4 + 2279 \alpha^3 \beta + 4269 \alpha^2 \beta^2 + 2279 \alpha \beta^3 + 274 \beta^4\\
%		&+ 225 \alpha^3 + 1025 \alpha^2 \beta + 1025 \alpha \beta^2 + 225 \beta^3\\
%		&+ 85 \alpha^2 + 205 \alpha \beta + 85 \beta^2\\
%		&+ 15 \alpha + 15 \beta + 1
	\end{split}
\end{equation*}
So the top Chern class is the degree $6$ term
\begin{equation*}
	\begin{split}
	c_6(\Sym^5Q)&=5\alpha(4\alpha+\beta)(3\alpha+2\beta)(2\alpha+3\beta)(\alpha+\beta)4\beta\\
	&=25\alpha\beta \left[4(\alpha+\beta)^2 + 9\alpha\beta\right]\left[6(\alpha+\beta)^2 + \alpha\beta\right]\\
	&=25\sigma_{1,1}(4\sigma_{1,0}^2+9\sigma_{1,1})(6\sigma_{1,0}^2+\sigma_{1,1})\\
	&=600\sigma_{1,1}\sigma_{1,0}^4+225\sigma_{1,1}^3+1450\sigma_{1,1}^2\sigma_{1,0}^2\\
	&=2875\sigma_{3,3}
	\end{split}
\end{equation*}

The last two equality made use of Tab. \ref{tab:multi}. Then we know the degree of top chern class $\on{deg}c_6(\Sym^5Q)=2875$, which is the answer to the problem raised at the front of this article, that the lines in a quintic 3-fold in $\bP^4$ is $2875$.

% TODO: add chern class of tensor
%
\sct*{Acknowledgement}
Thanks to Prof. James Fullwood, who guided and taught me during this directed study outside coursework by reading \tit{Enumerative geometry and string theory}\cite{Katz2006}. He helped me greatly both on study and life, I can never get a flavor of so many topics in modern algebraic geometry and mathematical physics at one time without this unforgettable experience.


% ----------------------------------------------------------------------------------------------------
\bibliographystyle{alpha}
\bibliography{../../universal-ref}
\end{document}
